\scanpage{8}
Finally, to express the holomorphy condition of the \unsure{fibration|filtration}, it is appropriate to introduce, \deleted{on} on the set \inserted{$\unsure{\mathscr{J}|\mathcal{J}|\gamma}$} of possible $j$ (or $j'$), a natural \uline{complex} \uline{analytic} \uline{manifold} structure, and the condition is expressed by the \uline{holomorphy} of the map
\[
s\longmapstoJ(s):\quad S\rightharpoonup\mathscr{J}.
\]

\deleted{In a \unsure{word}} Thus the set of \inserted{rigidified} $G$-subordinations with base $S$ is identified with the set $\operatorname{Mor}(S,\mathscr{J})$\margin{morphisms of analytic spaces}. Thus for variable $S$, one has a functor representable by $\mathscr{J}$.

\uline{NB} $G(\mathbb{R})$ acts in the obvious way on $\mathscr{J}$ (modulo inner automorphisms) ; the orbits are \uline{open}. Thus

\deleted{Let \deleted{\unsure{$G$ be a}} $V$ be a faithful $G$-module, and $V_0\subset V$ a finitely generated $Z$-module generating $V$. Then $\forall$ $G$-subordination on an $S$ defines a Hodge structure (\uline{not} modulo isogeny) on $V_0\times S$, Let $\Gamma\subset G(\mathbb{Q})$ be the stabilizer of $V_0$.}

the connected components of $\mathscr{J}$ are of the form $G(\mathbb{R})^{0}/H(\mathbb{R})\cap G(\mathbb{R})^{0}$, where $H$ is the centralizer of a $j:\unsure{\mathbb{S}|\mathrm{S}}\longrightarrow G$ in $G$ ; or again $G'(\mathbb{R})^{0}/H'(\mathbb{R})\cap G'(\mathbb{R})^{0}$, $H'$ the centralizer of a $j':\mathbb{U}\longrightarrow G'$ in $G'(\mathbb{R})^{0}$.
% archive sheet number: 7
